\section{Abstract}\label{sec:2-abstract}
Parameter estimation from pulsar timing array data is predominantly performed using the cross-power spectral density as the summary statistic fed into a likelihood or inference algorithm~\cite{renzini_stochastic_2022}.
While the power spectrum captures the mean spectral shape, it discards higher-order information encoded in the non-Gaussian features of the signal.
In cosmological large-scale structure analysis, morphological measures — specifically Minkowski functionals~\cite{mecke_robust_1994}, persistent homology~\cite{edelsbrunner_topological_2002}, and weighted morphological statistics~\cite{klatt_characterization_2022} — have been shown to carry significant additional constraining power beyond the power spectrum, particularly in the non-Gaussian regime~\cite{jalali_kanafi_imprint_2024,yip_cosmology_2024}.

This proposal investigates the application of these topological and geometrical summary statistics to SGWB data analysis in the context of PTA experiments.
The key observation is that the SGWB, as reconstructed from PTA cross-correlations across frequency bins and pulsar pairs, forms a two-dimensional map (pulsar pairs × frequency) whose morphological structure encodes information about the signal's spectral shape, angular distribution, and non-Gaussianity.
We propose to implement a systematic comparison of standard (power spectrum) and morphological summary statistics as inputs to a Simulation-Based Inference framework~\cite{cranmer_frontier_2020}, quantifying via Fisher forecasts and SBI posterior widths which combination of statistics yields the tightest constraints on cosmological source parameters.

The target sources include the cosmic string network and a simple power-law SGWB as a baseline.
This work establishes a direct bridge between Prof.
Movahed's morphological statistics toolkit~\cite{jalali_kanafi_imprint_2024,abedi_impact_2024} and gravitational wave data analysis, and is designed to yield a methods paper regardless of the physical source chosen.


\section{Objectives}\label{sec:2-objectives}
\begin{itemize}
    \item Implement Minkowski functionals~\cite{mecke_robust_1994} and persistent homology~\cite{edelsbrunner_topological_2002} as summary statistics operating on the frequency–pulsar-pair cross-correlation matrix of PTA residuals;
    \item Quantify via Fisher information matrix analysis the information gain of morphological statistics relative to the power spectrum alone for cosmic string parameter estimation~\cite{yip_cosmology_2024};
    \item Build a full SBI pipeline~\cite{cranmer_frontier_2020,alvey_simulationbased_2024} that uses these combined statistics as the compressed data representation fed to a neural ratio estimator;
    \item Identify the optimal combination of summary statistics as a function of signal-to-noise ratio and source model, producing a general result applicable to any SGWB source~\cite{renzini_stochastic_2022}.
\end{itemize}
\section{Methods}\label{sec:2-methods}

\begin{itemize}
    \item Mock PTA datasets are generated by injecting SGWB signals (cosmic string + SMBHB) into NANOGrav-15yr-calibrated noise realizations~\cite{afzal_nanograv_2023,johnson_nanograv_2024}, for a wide prior over source parameters.
    \item For each realization, the frequency-domain cross-correlation matrix C(f, a, b) (pulsars a, b) is computed. This matrix is the primary data object.
    \item Minkowski functionals (area, perimeter, Euler characteristic as a function of threshold)~\cite{mecke_robust_1994} are computed on excursion sets of C(f). Persistent homology~\cite{edelsbrunner_topological_2002} is computed using the Vietoris-Rips filtration of the frequency-pulsar graph.
    \item A Fisher information matrix forecast is constructed semi-analytically using the Gram matrix of summary statistic derivatives with respect to the source parameters~\cite{yip_cosmology_2024,jalali_kanafi_imprint_2024}.
    \item An SBI pipeline (TMNRE via \texttt{swyft}~\cite{miller_truncated_2021,miller_swyft_2022}) is constructed using each summary statistic set in turn, and posterior volumes are compared as a direct measurement of information content~\cite{alvey_simulationbased_2024}.
    \item The cluster is used for the simulation campaign ($\mathcal{O}(10^4$–$10^5)$ simulations) and for parallelized morphological statistic computation.
\end{itemize}